\chapter*{A Note on Consequences: What Misalignment, Short of Catastrophe, Looks Like}
\addcontentsline{toc}{chapter}{A Note on Consequences}
\markboth{A Note on Consequences}{}

The book has built one argument and one apparatus. The argument is that the specification problem is structural and that capability amplifies it. The apparatus is the mathematics of optimal agency plus the mitigations that push back the threshold at which catastrophic failure becomes hard to detect.

That argument and apparatus naturally lead to a register of stakes the book treated: the existential register. Misaligned ASI taking control. Sufficiently capable optimizers of misspecified objectives producing catastrophic outcomes. The math points there because the math is about optimization, and the worst case for an optimizer with a wrong objective is, formally, the worst case.

There is a second register of stakes that the book did not develop in the main chapters, because doing so would have required a different kind of book. The second register is what happens to people and societies when the gap is not closed, in outcomes short of full catastrophe. These are the more probable near-term outcomes. They are also the outcomes most readers will live in.

This note is a coda. It is brief. The book's apparatus illuminates these consequences without committing to specific forecasts; the thinking about the consequences is the reader's to do.

\section*{The resource-curse parallel}

States whose GDP comes largely from resources rather than from human productivity tend to underinvest in their populations. The historical pattern is well documented. Oil states, mineral states, rentier states: Saudi Arabia, Venezuela, Equatorial Guinea, Angola, the Gulf monarchies. The mechanism is structural. If the population is not the source of value, the state's incentive to invest in education, infrastructure, healthcare, and welfare is weaker than in states where the population is what creates wealth. The labor of citizens is not what pays for the state's existence. The resource is.

This is the resource curse. It is not deterministic, but it is the empirical default. The historical analogues are clear enough that political economists have a vocabulary for them.

The relevant question for the book's reader is what happens when AGI or ASI provides most of the economic value of a civilization, and human labor is largely surplus to that value.

The same incentive structure applies, at species or civilizational scale. The ``country of geniuses in a data center'' framing made popular by Dario Amodei is exactly this. A small group with concentrated capability, in the form of advanced AI systems, produces economic output that does not require most of the human population as input. The political economy then resembles a rentier state, not a productive economy. The population's labor is no longer the source of value. The AI is.

The historical record of rentier states is not a model the reader will find encouraging. Under the resource-curse incentive, the state's investment in human capability declines, because human capability is no longer load-bearing. Public education, healthcare, civic infrastructure, democratic participation, the welfare state in its various forms: each was built on the assumption that the population is the source of productive value. Remove that assumption and the institutions face structural pressure that no historical analogue cleanly handles.

This is not a forecast. It is the application of a known incentive structure to a new technological condition. Whether it plays out this way depends on what societies choose to do, what political pressure their populations can exert, what alternative arrangements they construct. The book's apparatus does not predict the choices; it points at what the incentives default to in the absence of countervailing forces.

\section*{Identity and skill}

The second consequence concerns identity rather than political economy.

Humans currently locate self-worth in capability. What we can do. What we can build. What we can understand. What we can make. The skills we have are partly how we identify ourselves to ourselves. What we are good at is partly what we mean when we say we are someone.

There is a particular cruelty in the order this arrives in, and Chapter 16's jagged edge named it. The twentieth century taught us to expect automation from the bottom up: machines take the manual work first, the routine before the skilled, and human dignity retreats upward into thought. This wave runs the other way. The competences evolution spent the most optimization on, a steady pair of hands, a read of a room, balance on a wet floor, are the last to fall, because they are the hardest to specify. The abstract, verbal, analytical work we built our most prestigious identities around is comparatively easy for these systems, and it is going first. The lawyer and the radiologist are exposed before the plumber and the nurse. The transition does not come for the work we were taught to consider beneath us. It comes for the work we were taught to be proud of.

AI better at most things is a different world for identity. There are two cases.

The pessimistic case is meaning collapse at scale. Industrial revolutions produced versions of this when artisanal identities were displaced by factories. Automation produced versions when factory work was displaced by services. Each transition was real and painful even when partial, and each took decades to settle. The transition the book's apparatus points at is faster and broader. The historical analogues do not promise it will settle gracefully.

The optimistic case is worth giving space to, because it has more empirical evidence on its side than the pessimistic case has on its.

\section*{The twelve-year-old-girls observation}

Humans still play chess against other humans even though Stockfish has been crushing humans for decades. We still race each other even though cheetahs beat us and cars beat the cheetahs. We still write poems and read poems even though large language models can produce text that scans, scans well, and is correctly metered. The story that matters at human scale, in human time, is a human story. The presence of a higher tier of capability does not, by itself, evacuate the human-scale story of meaning.

The clearest illustration of this is sociological rather than competitive. Watch a gaggle of twelve-year-old girls in any setting. They take an intense interest in each other's lives. They care, in great detail, about who said what to whom, about who is friends with whom, about what was worn, about what was overheard, about whose phone was buzzing, about what someone's mother said at the door. The authoritative adults who have nominal dominion over their lives are visible in the background, looming, but their pronouncements register at most as inflection in the larger texture of what is going on at twelve-year-old scale.

The Charlie Brown image compresses this. In the Peanuts cartoons, the adults speak in a noise that is canonically rendered as \emph{wah wah wah}. The children's world goes on. The adults are real, are authoritative, are not contested. They are also not where the meaning is. The meaning is in what Lucy said to Linus, what Snoopy is doing on the doghouse, what Charlie Brown's misery is this week. The higher tier is real and is not the location of the story.

Humans, as social animals, deeply care about what other humans think, for a variety of reasons both good and bad. This is sometimes called speciesism, and as a moral position it has its critics, but as a descriptive fact about human attention it is well-attested. We attend to other humans. We attend to them when we are children and adults loom. We attend to them when we are adults and other tiers (governments, corporations, gods, fates) loom. The attention to each other is not contingent on being the most capable thing in the environment.

Whether this human-to-human attention will survive a truly general agentic intelligence acting as a countervailing force is an open question. The book cannot answer it. What can be observed is that the relevant comparison cases (a higher tier in the environment that is not the location of human-scale meaning) have been compatible with rich, dense, meaningful human-scale life. The twelve-year-old girls are not paying attention to the adults. They have plenty of meaning anyway.

This is the optimistic case. It is not a guarantee. It is a noticing.

\section*{The non-catastrophic-but-bad case}

The two registers, resource curse and identity, compose.

AGI controlled by a small group, without losing control to misalignment, is not existential risk. It is the application of the resource-curse incentive at species scale. The elite class faces reduced reason to invest in the rest, because the rest is not the source of value. Most current institutions were built on the assumption that the population is the source of productive value. Remove that assumption and the institutions face pressure.

This is a bad outcome. It is not as bad as existential catastrophe. Whether it is more or less probable than catastrophe is a question the book does not try to answer; both possibilities follow from the same structural fact, and the relative weights one assigns are a matter of judgment the math does not settle.

The book did not develop this thread in the main chapters because the main argument is about the technical gap. The technical gap is what makes catastrophic outcomes possible. The non-catastrophic outcomes are what happens when the gap is partially open: capability concentrated in a few systems, controlled by a few people, without the broad alignment that would distribute its benefits to the population.

The countervailing forces are political, not technical. They are about how societies arrange the ownership and governance of these systems. Those questions are real and important and outside what the book set out to do. They are also questions the reader who has the book's apparatus is well-positioned to think about, because the apparatus names what is structurally at stake. The political work follows.

\section*{Closing}

The book gave you a way of thinking about the technical gap between AIXI and the systems being built now. The stakes of that gap split into two registers. The existential register, which the math points at directly, the book treated in the main chapters. The broad register, which the math points at indirectly, this note has gestured at.

The reader has the apparatus. The thinking is the reader's to do. The book's argument that capability without alignment is the central technical problem of our era is one input to that thinking. The reader's own observations of what kinds of communities, institutions, and lives they want to be part of are another. The combination is what produces the answer to the question the book did not answer: what should we do.

The book is what it is. The thinking is yours.
